\documentclass[11pt]{article}

\usepackage[a4paper,margin=2.5cm]{geometry}
\usepackage{amsmath,amssymb,amsthm}
\usepackage[colorlinks=true,linkcolor=blue,citecolor=blue,urlcolor=blue]{hyperref}
\usepackage{booktabs}
\usepackage{enumitem}
\usepackage{natbib}

\newtheorem{theorem}{Theorem}[section]
\newtheorem{lemma}{Lemma}[section]
\newtheorem{proposition}{Proposition}[section]
\newtheorem{corollary}{Corollary}[section]
\newtheorem{definition}{Definition}[section]
\newtheorem{remark}{Remark}[section]

\newcommand{\twoI}{2I}
\newcommand{\Vsixhundred}{V_{600}}
\newcommand{\Dicfive}{\mathrm{Dic}_5}
\newcommand{\Vtwentyfour}{V_{24}}
\newcommand{\sigmaG}{\sigma}
\newcommand{\tausigma}{\tau_{\!\sigmaG}}
\newcommand{\tausigmaj}[1]{\tau^{(#1)}_{\!\sigmaG}}
\newcommand{\Ttau}{T_{\tau}}
\newcommand{\Trbil}{\mathrm{Tr}}

\title{The $\tausigma$ Involution: A Canonical Galois-Coset Duality\\
       on the Binary Icosahedral Group}
\author{%
  Lee Smart \\[3pt]
  \textit{Institute of Vibrational Field Dynamics} \\[2pt]
  \texttt{contact@vibrationalfielddynamics.org}}
\date{May 2026}

\begin{document}

\maketitle

\noindent\textbf{Status:} Preprint (not peer-reviewed). Pure mathematics; no physical interpretation is claimed.

\begin{abstract}
We construct an explicit involution
$\tausigma : \Vsixhundred \to \Vsixhundred$ on the vertex set of the
$600$-cell, equivalently the binary icosahedral group $\twoI$. The
construction is a vertex permutation rather than a Euclidean
reflection; testing two restricted families fixing the binary
dihedral subgroup $H = \Dicfive$ pointwise — single
$E_8$-Weyl-reflection candidates (which act trivially on
$\Vsixhundred$ in every tested case) and parity-flip candidates
(which produce zero closed-and-involution target hits) — motivates
using a vertex permutation instead (Remark~\ref{rem:reflfail}; this is
a computational observation, not an exhaustive non-existence claim,
and W(H$_4$) was not separately enumerated). The construction instead
uses the right-coset
$\Ttau$-cycle structure paired with a $\sigmaG$-Galois projection in
the icosian trace metric. For each non-bulk right coset
$Hg \subset \twoI$, which contains exactly one whole $\Ttau$-cycle of
K-class $52$ and one whole $\Ttau$-cycle of K-class $20$ (Paper~1
\citep{V600Paper1}), the construction selects a $\Ttau$-equivariant
phase swap by the symmetric-maximin criterion on the trace inner
product $\Trbil_{K/\mathbb{Q}}\,\mathrm{Re}(\sigmaG(v) \cdot \bar w)$.
The phase selection yields exactly two admissible phases per coset
(the antipodal pair $j \in \{0, 5\}$), giving a $\mathbb{Z}_2^5 = 32$-fold
canonical lift. We prove $\tausigma$ satisfies four canonical
properties: it is an involution, fixes $H$ pointwise, exchanges the
$K=52$ and $K=20$ cycle classes within each non-bulk right coset, and
commutes with $\Ttau$ (T-equivariance). It additionally commutes with
the antipodal map $v \mapsto -v$. The canonical phase
$(0,0,0,0,0)$-lift is exhibited explicitly as a $120$-vertex map and
verified by an exact-arithmetic certificate.
\end{abstract}

\section{Introduction}\label{sec:intro}

The binary icosahedral group $\twoI$ — equivalently the vertex set
$\Vsixhundred$ of the regular $600$-cell — admits a rich finite-group
structure that has been studied for over a century in the contexts of
regular polytopes \citep{Coxeter1973}, lattice theory
\citep{ConwaySloane1988,Elkies1995}, and Coxeter / reflection groups
\citep{Humphreys1990}. Recent work \citep{V600Paper1} has identified a
distinguished subgroup chain $H = \Dicfive \subset \twoI$ ($|H| = 20$)
arising from the bulk decomposition of $\Vsixhundred$ under
$\Ttau$-cycle K-class structure, with $H$ self-normalizing in
$\twoI$.

In this note we construct an involution
$\tausigma : \Vsixhundred \to \Vsixhundred$ that

\begin{enumerate}[leftmargin=2em]
\item fixes $H$ pointwise,
\item swaps the $K=52$ and $K=20$ cycle classes within each non-bulk
right coset of $H$,
\item commutes with $\Ttau$, and
\item commutes with the antipodal map.
\end{enumerate}

The motivation is internal to the V$_{600}$ programme: a subsequent
paper uses this involution to state phase-independence for
operator-trace identities on the $\Ttau$-cycle space of
$\Vsixhundred$. The construction is purely combinatorial-arithmetic;
no physical interpretation is claimed in the present paper.

A natural first attempt — a Euclidean reflection of $\Vsixhundred$
fixing $H$ — was tested in two restricted families:
\texttt{block3c\_e8\_reflection\_obstruction.py} reports that
single-Euclidean-reflection candidates fixing $H$ act trivially on
$\Vsixhundred$ in the tested embedding, and
\texttt{block3b\_parity\_flip\_tau\_sigma.py} reports zero
target hits in the parity-flip family. We summarise both in
Section~\ref{sec:reflfail} (no exhaustive non-existence claim is made,
and W(H$_4$) was not separately enumerated), then turn to a
vertex-permutation construction in Section~\ref{sec:construction}. Section~\ref{sec:main}
proves the four canonical properties of the resulting map; Section~\ref{sec:Z25}
discusses the $\mathbb{Z}_2^5$ phase ambiguity.

\section{Setup}\label{sec:setup}

We import the following from Paper 1 \citep{V600Paper1}; the present
paper does not re-prove any of these statements.

\begin{itemize}[leftmargin=2em]
\item $K := \mathbb{Q}(\sqrt{5})$ with Galois twist
$\sigmaG : a + b\sqrt{5} \mapsto a - b\sqrt{5}$, extended componentwise
to the icosian quaternion algebra $\mathbb{H}_K$.
\item $\Vsixhundred = \twoI \subset \mathrm{Sp}(1, K)$: the $120$ unit
icosian quaternions forming the binary icosahedral group, partitioned
into $8$ axis, $16$ half-type, and $96$ even-permutation vertices.
\item $\Vtwentyfour := \Vsixhundred \cap \mathrm{Fix}(\sigmaG)$: the
24-cell sublattice ($8$ axis + $16$ half-type vertices).
\item $\Ttau : g \mapsto \tau g$ for some chosen $\tau \in \twoI$ with
$\mathrm{ord}(\tau) = 10$. The cycle decomposition of $\Ttau$
partitions $\Vsixhundred$ into 12 cycles each of length 10, with
K-multiset $\{72{:}1,\,0{:}1,\,52{:}5,\,20{:}5\}$.
\item $H := \Dicfive$, the bulk subgroup of order $20$ formed by the
union of the $K=72$ and $K=0$ cycles. Self-normalizing and non-normal
in $\twoI$ \citep{V600Paper1}.
\end{itemize}

\begin{lemma}[Right-coset carrier; Paper 1, Coset cycle composition]\label{lem:rightcoset}
For every non-bulk right coset $Hg \subset \twoI$, the coset is
exactly the disjoint union of one whole $\Ttau$-cycle of K-class $52$
and one whole $\Ttau$-cycle of K-class $20$. In particular $\Ttau$
preserves $Hg$ setwise (since $\tau \in H$, where $\tau$ is the chosen
order-$10$ generator of the bulk $\Ttau$-cycle of $K=72$ that lies
inside $H = \Dicfive$, we have $\tau H = H$ and hence
$\Ttau(Hg) = (\tau H) g = Hg$); the K-class pairing $52$ with
$20$ is the additional content of the lemma and is established in
Paper 1 by direct enumeration. The five non-bulk right cosets
exhaust the ten non-bulk $\Ttau$-cycles, paired one $K=52$ with one
$K=20$.
\end{lemma}

By contrast, non-bulk \emph{left} cosets $gH$ contain two vertices from
each of the ten non-bulk $\Ttau$-cycles and are \emph{not} whole-cycle
carriers (Paper 1, ``Coset cycle composition'' lemma). For the construction below it is
essential to use the right-coset family.

\subsection{Trace inner product}

The icosian quaternion algebra $\mathbb{H}_K$ with $K = \mathbb{Q}(\sqrt{5})$
carries a natural rational trace bilinear form,
\[
  \langle v, w \rangle_{\mathrm{tr}}
  \;:=\;
  \mathrm{Tr}_{K/\mathbb{Q}}\!\left(\,\mathrm{Re}(v \cdot \bar w)\,\right),
\]
where $\bar w$ is the quaternionic conjugate, $\mathrm{Re}$ extracts the
real (scalar) component, and $\mathrm{Tr}_{K/\mathbb{Q}}(a + b\sqrt{5}) = 2a$.
Up to a global rescaling, this form puts $\Vsixhundred$ inside the
$E_8$ lattice \citep{Elkies1995,ConwaySloane1988}; on the cross-Galois
pairs used in the construction below the values are rational (and in
fact lie in $\frac{1}{4}\mathbb{Z}$), and all arithmetic is exact in
$\mathbb{Q}(\sqrt{5})$.

\section{Why Euclidean reflections fail}\label{sec:reflfail}

The natural first attempt at constructing $\tausigma$ is a Euclidean
reflection of $\Vsixhundred$ fixing $H$ pointwise. The following
computational observation motivates abandoning this route.

\begin{remark}[Reflection obstruction, computational]\label{rem:reflfail}
Two auxiliary script families in the V$_{600}$ programme repository
test isometric $\tausigma$ candidates that fix $H = \Dicfive$
pointwise:
\begin{itemize}[leftmargin=2em]
\item \texttt{block3c\_e8\_reflection\_obstruction.py} (with the
embedding helpers in \texttt{block3a\_e8\_icosian\_lattice.py}) checks
single Euclidean-reflection candidates fixing $H$ and reports that
each such candidate acts trivially on $\Vsixhundred$ in the tested
embedding;
\item \texttt{block3b\_parity\_flip\_tau\_sigma.py} searches the
parity-flip family for a $\Vsixhundred$-permutation that fixes $H$
pointwise, swaps the $K=52$ and $K=20$ cycle classes, and is closed
under $\twoI$, and reports zero target hits.
\end{itemize}
No exhaustive Weyl-group enumeration is claimed, and we do not assert
a closed-form non-existence theorem; this remark only motivates the
vertex-permutation construction below.
\end{remark}

The remainder of the paper does not depend on Remark~\ref{rem:reflfail}
mathematically; it is included only to motivate why we construct
$\tausigma$ as a vertex permutation rather than as a reflection. The
construction proceeds via a $\sigmaG$-Galois projection on right
cosets.

\section{The $\sigmaG$-projection construction}\label{sec:construction}

\subsection{Cycle phase swap}

For a non-bulk right coset $Hg \subset \twoI$, fix base points
$v_0 \in C_{52}$ and $w_0 \in C_{20}$ (both supplied by the
deterministic vertex enumeration of
\texttt{vfd\_v600.icosian.build\_vertices()}; see Section~\ref{sec:Z25}
for the convention dependence) and write the constituent
$\Ttau$-cycles as
\[
  C_{52} = (v_0, v_1, \ldots, v_9), \qquad
  C_{20} = (w_0, w_1, \ldots, w_9),
\]
with $v_{k+1} = \tau v_k$ and $w_{k+1} = \tau w_k$ (indices mod $10$).
The phase labels below are computed relative to this choice; a
different pair $(v_0, w_0)$ would relabel them.

\begin{definition}[Phase-$j$ swap]\label{def:phaseswap}
For each $j \in \mathbb{Z}/10$, the \emph{phase-$j$ swap} is the
involution $\tausigmaj{j} : Hg \to Hg$ defined by
\[
  \tausigmaj{j}(v_k) := w_{(j + k) \bmod 10}, \qquad
  \tausigmaj{j}(w_{(j + k) \bmod 10}) := v_k.
\]
Each $\tausigmaj{j}$ exchanges the two cycles, is a well-defined
involution on $Hg$, and is $\Ttau$-equivariant by construction.
\end{definition}

There are $10$ phases per coset, so naively $10^5 = 10^5$ candidates
across the five non-bulk cosets. The $\sigmaG$-projection criterion
below restricts to $2^5 = 32$ canonical lifts.

\subsection{Symmetric-maximin $\sigmaG$-projection}

\begin{definition}[Symmetric maximin score]\label{def:maximin}
For a non-bulk right coset $Hg$ with cycle data
$(C_{52}, C_{20}) = (v_k, w_k)$, define the forward and backward
$\sigmaG$-projections
\[
  s_{\to}(j,k) \;:=\; \langle \sigmaG(v_k),\; w_{(j+k) \bmod 10} \rangle_{\mathrm{tr}},
  \qquad
  s_{\leftarrow}(j,k) \;:=\; \langle \sigmaG(w_{(j+k) \bmod 10}),\; v_k \rangle_{\mathrm{tr}},
\]
and the \emph{symmetric phase score}
\[
  s(j) \;:=\; \min_{k \in \mathbb{Z}/10}\,
  \min\bigl(s_{\to}(j,k),\, s_{\leftarrow}(j,k)\bigr).
\]
The \emph{maximin admissible phase set} is
\[
  \Phi(Hg) \;:=\; \mathrm{argmax}_{j \in \mathbb{Z}/10}\, s(j).
\]
The two-sided minimum makes $s(j)$ symmetric in the role of the two
cycles and is the form verified by the companion script
\texttt{verify.py}.
\end{definition}

\begin{theorem}[Phase selection]\label{thm:phases}
Under the base-point convention of Section~\ref{sec:construction}
(both $v_0 \in C_{52}$ and $w_0 \in C_{20}$ supplied by the
deterministic ordering of
\texttt{vfd\_v600.icosian.build\_vertices()}), for every non-bulk
right coset $Hg \subset \twoI$, $\Phi(Hg) = \{0,\, 5\}$, with
$s(0) = s(5) = 0$ and $s(j) = -\frac{5}{4}$ for
$j \in \{1,2,3,4,6,7,8,9\}$. The two winning phases are separated by
$5$ (i.e., antipodal indices on the length-$10$ cycle); the labels
$\{0, 5\}$ themselves are convention-dependent
(Remark~\ref{rem:conv}). Phase $j=5$ is the antipodal shift of $j=0$.
\end{theorem}

\begin{proof}
Direct exact-arithmetic computation in $\mathbb{Q}(\sqrt{5})$ (the
companion script \texttt{verify.py} enumerates all five non-bulk right
cosets and all $10$ phases, asserts the full symmetric score table
$s(0) = s(5) = 0$ and $s(j) = -\tfrac{5}{4}$ otherwise, and hence
verifies the argmax). The identity
$j=5 \leftrightarrow j=0 + \mathrm{antipode}$ follows from the central
involution $\tau^5 = -1 \in \twoI$ acting on cycle indices by
$k \mapsto k+5 \bmod 10$.
\end{proof}

\subsection{Canonical lift}

Definition~\ref{def:phaseswap} and Theorem~\ref{thm:phases} together
yield $|\Phi(Hg)|^5 = 2^5 = 32$ admissible compositions across the five
non-bulk right cosets. We name the all-zero choice canonical.

\subsection{Algorithm (pseudocode)}

The construction and verification can be summarised as the following
explicit pseudocode (the script \texttt{verify.py} is a literal
translation):

\begin{verbatim}
build state via vfd_v600.group.build_state()
   -> verts, cycles, K_of_cycle, Dic5, T_tau

for each non-bulk right coset Hg of H = Dic_5:
    let C52 = the K=52 cycle in Hg, base point v_0
    let C20 = the K=20 cycle in Hg, base point w_0
    for j in 0..9:
        s(j) := min over k in 0..9 of
                 min( <sigma(v_k), w_{(j+k) mod 10}>_tr,
                      <sigma(w_{(j+k) mod 10}), v_k>_tr )
    Phi(Hg) := argmax_{j} s(j)         # equals {0, 5}

choose phase j_g = 0 for each Hg (canonical lift)
build tau_sigma:
    tau_sigma(h)        := h            for h in Dic_5
    tau_sigma(v_k)      := w_{(j_g+k) mod 10}
    tau_sigma(w_{j_g+k}):= v_k

verify: P1 (involution), P2 (fix Dic_5),
        P3 (K=52 <-> K=20), P4 (T_tau-equivariance),
        antipodal: tau_sigma(-v) == -tau_sigma(v),
        right-coset preservation,
        match against data/tau_sigma_canonical.txt.
\end{verbatim}

\begin{definition}[Canonical $\tausigma$]\label{def:canonical}
The \emph{canonical $\tausigma$} is the involution
$\tausigma : \Vsixhundred \to \Vsixhundred$ obtained by:
\begin{itemize}[leftmargin=2em]
\item $\tausigma(h) = h$ for every $h \in H$;
\item For each of the five non-bulk right cosets $Hg$, applying the
phase-$0$ swap $\tausigmaj{0}$ of Definition~\ref{def:phaseswap}.
\end{itemize}
The 120-vertex map is recorded explicitly in
\texttt{papers/v600-programme/data/tau\_sigma\_canonical.txt}.
\end{definition}

\subsection{Worked example}

For the canonical state with $\tau$ chosen as the first
order-$10$ icosian in the deterministic vertex enumeration, the first
non-bulk right coset has $K=52$ cycle and $K=20$ cycle
\[
  C_{52} = [4, 119, 81, 80, 118, 5, 112, 86, 87, 113],\qquad
  C_{20} = [6, 93, 59, 58, 92, 7, 90, 60, 61, 91].
\]
The phase-$0$ swap pairs them index-wise:
$4 \leftrightarrow 6$, $119 \leftrightarrow 93$, $81 \leftrightarrow 59$,
$80 \leftrightarrow 58$, $118 \leftrightarrow 92$, $5 \leftrightarrow 7$,
$112 \leftrightarrow 90$, $86 \leftrightarrow 60$, $87 \leftrightarrow 61$,
$113 \leftrightarrow 91$.

\section{Main theorem}\label{sec:main}

\begin{theorem}[Canonical $\tausigma$ properties]\label{thm:main}
The canonical $\tausigma$ of Definition~\ref{def:canonical} satisfies:
\begin{enumerate}[label=(P\arabic*),leftmargin=3em]
\item \emph{Involution:} $\tausigma \circ \tausigma = \mathrm{id}_{\Vsixhundred}$.
\item \emph{Fixes the bulk subgroup:} $\tausigma(h) = h$ for every $h \in H$.
\item \emph{Cross-K coset swap:} for every $v \in Hg$ with $Hg$
non-bulk, the K-class of the cycle of $\tausigma(v)$ is the
K-opposite of the K-class of the cycle of $v$ (i.e., $52
\leftrightarrow 20$).
\item \emph{$\Ttau$-equivariance:} $\tausigma \circ \Ttau = \Ttau \circ \tausigma$.
\end{enumerate}
Additionally, $\tausigma$ commutes with the antipodal map
$v \mapsto -v$.
\end{theorem}

\begin{proof}
(P1) Each phase-$0$ swap on a non-bulk right coset is an involution by
Definition~\ref{def:phaseswap}, and identity is an involution on $H$;
the disjoint union of involutions on the six cosets is an involution
on $\Vsixhundred$.

(P2) Direct from the definition.

(P3) Each phase-$0$ swap exchanges $C_{52}$ and $C_{20}$ in its right
coset; the K-class assignment is $\Ttau$-cycle-determined, so each
swap pair has K-classes $\{52, 20\}$.

(P4) For $v = v_k \in C_{52}$ in some non-bulk right coset,
$\tausigma(\Ttau(v_k)) = \tausigma(v_{k+1}) = w_{k+1} =
\Ttau(w_k) = \Ttau(\tausigma(v_k))$, with indices mod 10. Same for
$v = w_k$. For $v \in H$, $\tausigma(v) = v$ and $\Ttau(v) =
\tau v \in H$ also fixed.

Antipodal compatibility. The element $\tau \in \twoI$ has order $10$
in the binary icosahedral group; the unique central involution of
$\twoI$ is $-1$, and the order-$10$ subgroup $\langle \tau \rangle$
contains it as $\tau^5 = -1$ (this is the standard unit-quaternion
fact $|\tau| = 1$ with $\tau^{10} = 1$ forcing $\tau^5 \in \{\pm 1\}$;
$\tau^5 = 1$ would contradict $\mathrm{ord}(\tau) = 10$). Hence the
antipodal map acts on cycle indices by
$v_k \mapsto -v_k = \tau^5 v_k = v_{k+5}$
(and similarly $w_k \mapsto w_{k+5}$), with indices mod $10$. Both
$v_{k+5}$ and $w_{k+5}$ remain in the same right coset $Hg$, since
right cosets are setwise $\Ttau$-invariant by
Lemma~\ref{lem:rightcoset}. Therefore
\[
  \tausigma(-v_k) \;=\; \tausigma(v_{k+5}) \;=\; w_{k+5} \;=\; -w_k \;=\; -\tausigma(v_k),
\]
and analogously for $v = w_k$. For $v \in H$, $\tausigma(v) = v$ and
$-v \in H$ (since $-1 \in H$), so $\tausigma(-v) = -v = -\tausigma(v)$.
\end{proof}

\section{The $\mathbb{Z}_2^5$ phase ambiguity}\label{sec:Z25}

By Theorem~\ref{thm:phases}, each non-bulk right coset admits two
maximin phases, $\{0, 5\}$. Choosing a phase independently for each
of the five non-bulk right cosets gives $2^5 = 32$ admissible
involutions, all satisfying (P1)–(P4) of Theorem~\ref{thm:main}. The
group of phase choices is $\mathbb{Z}_2^5$, generated by the five
per-coset antipodal shifts $j \mapsto j + 5$.

\begin{remark}[Convention vs.~freedom]\label{rem:conv}
The labels $\{0, 5\}$ for the maximin phase set are
\emph{convention-dependent}; what is intrinsic is that
$|\Phi(Hg)| = 2$ and that the two winning phases are antipodal
(separated by $5$ on the length-$10$ cycle, equivalently related by
$\tau^5 = -1$). The specific labels arise from:
(i) the choice of base points $v_0 \in C_{52}$ and $w_0 \in C_{20}$ in
Definition~\ref{def:phaseswap}; and
(ii) the deterministic ordering of vertices supplied by
\texttt{vfd\_v600.icosian.build\_vertices()}. A different choice of
base points relabels the phases by a shift $j \mapsto j + c$ for some
$c \in \mathbb{Z}/10$, but the $\mathbb{Z}_2^5$ orbit of canonical
lifts is intrinsic to the construction.
\end{remark}

The certificate \texttt{verify.py} enumerates all $32$ phase
combinations and confirms that every one yields an admissible
$\tausigma$ involution satisfying (P1)–(P4).

\section{Conclusion}\label{sec:concl}

The binary icosahedral group $\twoI$ admits a canonical $\sigmaG$-Galois
involution $\tausigma$ satisfying four structural properties (P1)–(P4)
of Theorem~\ref{thm:main}, plus antipodal compatibility. The
construction is a vertex permutation (motivated, but not required, by
the computational reflection obstruction recorded in
Remark~\ref{rem:reflfail}), realised via cycle-phase swaps on right
cosets of the bulk subgroup $\Dicfive$ selected by symmetric-maximin
$\sigmaG$-projection in the icosian trace metric
(Theorem~\ref{thm:phases}). The $\mathbb{Z}_2^5$
ambiguity in phase choice yields $32$ canonical lifts; we name the
all-zero choice the canonical representative and exhibit its
120-vertex map explicitly.

\appendix

\section{Verification certificate}\label{app:verify}

The companion script \texttt{verify.py} performs the following exact
checks:
\begin{enumerate}[leftmargin=2em]
\item $|\Vsixhundred| = 120$, $|H| = 20$, $|\Vtwentyfour| = 24$.
\item Every non-bulk right coset $Hg$ is one whole $\Ttau$-cycle of
$K=52$ plus one whole $\Ttau$-cycle of $K=20$.
\item Regression: every non-bulk left coset $gH$ contains exactly
two vertices from each of the ten non-bulk $\Ttau$-cycles
(\emph{not} a whole-cycle carrier).
\item Per non-bulk right coset, the symmetric maximin score $s(j)$ of
Definition~\ref{def:maximin} satisfies $s(0) = s(5) = 0$ and
$s(j) = -\tfrac{5}{4}$ for $j \in \{1,2,3,4,6,7,8,9\}$, hence
$\Phi(Hg) = \{0, 5\}$.
\item The phase-$(0,0,0,0,0)$ rebuild matches the loaded canonical
map \texttt{data/tau\_sigma\_canonical.txt} vertex-by-vertex.
\item The rebuilt map satisfies (P1)–(P4) and antipodal compatibility.
\item Right-coset preservation: $\tausigma$ preserves each right coset
$Hg$ setwise (post-build check, not part of the theorem statement; it
holds by Definition~\ref{def:canonical}, since each per-coset phase
swap is internal to its right coset, and is checked separately by the
script — properties (P1)--(P4) alone do not imply it, since one could
in principle pair a $K=52$ cycle of one coset with a $K=20$ cycle of
another and still satisfy (P1)--(P4)).
\item All $2^5 = 32$ phase combinations yield valid involutions
satisfying (P1)–(P4).
\end{enumerate}
The script depends only on the Python standard library and the shared
\texttt{vfd\_v600} package. Reproducibility command:
\begin{verbatim}
PYTHONPATH=papers/v600-programme/lib \
    python3 papers/tau-sigma-construction/verify.py
\end{verbatim}

\bibliographystyle{plainnat}
\bibliography{references}

\end{document}
